\documentclass[aspectratio=169, 10pt, spanish]{beamer}

\usepackage[utf8]{inputenc}
\usepackage[T1]{fontenc}
\usepackage{lmodern}
\usepackage{helvet}
\renewcommand{\familydefault}{\sfdefault}
\usepackage[spanish, es-tabla, es-nodecimaldot]{babel}
\usepackage{tikz}
\usetikzlibrary{shapes.geometric, arrows.meta, positioning, calc, fit, backgrounds}
\usepackage{tcolorbox}
\tcbuselibrary{skins, breakable}
\usepackage{amsmath, amssymb}
\usepackage{colortbl}
\usepackage{booktabs}
\usepackage{tabularx}

% Configuración general de Beamer
\setbeamertemplate{navigation symbols}{}
\setbeamercolor{background canvas}{bg=white}

% Paleta de Colores Blueprint (Cátedra Msc. Raúl Huillca)
\definecolor{bpNavy}{RGB}{28, 43, 54}        % Barra superior, títulos y marcos
\definecolor{bpBlue}{RGB}{43, 91, 132}       % Proceso [P] / Propuesta / Tecnología
\definecolor{bpGold}{RGB}{200, 157, 66}      % Tecnología [T] / Núcleo / Método
\definecolor{bpRust}{RGB}{179, 92, 62}       % Contexto [C] / Ámbito / Aplicación
\definecolor{bpDark}{RGB}{33, 37, 41}        % Texto principal
\definecolor{bpLightGray}{RGB}{245, 247, 250}% Fondo técnico
\definecolor{bpGridLine}{RGB}{225, 230, 238} % Color cuadrícula blueprint
\definecolor{bpGreen}{RGB}{46, 125, 50}      % Validado / Aprobado / Fase 3
\definecolor{bpRed}{RGB}{198, 40, 40}        % Error / Anti-patrón

% Personalización de Frametitle estilo Blueprint
\setbeamertemplate{frametitle}{
    \begin{tikzpicture}[remember picture, overlay]
        \fill[bpNavy] (current page.north west) rectangle ([yshift=-1.05cm]current page.north east);
        \node[anchor=west, text=white, font=\bfseries\sffamily\normalsize] at ([xshift=0.6cm, yshift=-0.52cm]current page.north west) {
            \insertframetitle
        };
        \node[anchor=east, text=white!75!white, font=\tiny\sffamily] at ([xshift=-0.6cm, yshift=-0.52cm]current page.north east) {
            SEMINARIO DE INVESTIGACIÓN II
        };
    \end{tikzpicture}
    \vspace{0.25cm}
}

% Pie de página estilo Plano
\setbeamertemplate{footline}{
    \begin{tikzpicture}[remember picture, overlay]
        \draw[bpGridLine, line width=0.8pt] ([yshift=0.5cm]current page.south west) -- ([yshift=0.5cm]current page.south east);
        \node[anchor=west, text=bpNavy!80!black, font=\tiny\sffamily] at ([xshift=0.6cm, yshift=0.25cm]current page.south west) {
            UNSAAC -- Ingeniería Informática y de Sistemas $\mid$ Br. Edmil Jampier Saire Bustamante
        };
        \node[anchor=east, text=bpNavy!80!black, font=\tiny\sffamily] at ([xshift=-0.6cm, yshift=0.25cm]current page.south east) {
            \insertframenumber{} / \inserttotalframenumber
        };
    \end{tikzpicture}
}

\begin{document}

% ==========================================================
% SLIDE 1: PORTADA TÉCNICA (BLUEPRINT)
% ==========================================================
\begin{frame}[plain]
    \begin{tikzpicture}[remember picture, overlay]
        % Cuadrícula blueprint tenue
        \draw[step=0.5cm, bpGridLine, line width=0.3pt] (current page.south west) grid (current page.north east);
        
        % Marco técnico ajustado
        \draw[bpNavy, line width=1.4pt] ([xshift=0.5cm, yshift=0.5cm]current page.south west) rectangle ([xshift=-0.5cm, yshift=-0.5cm]current page.north east);
        
        % Marcas de cota / miras en esquinas
        \draw[bpNavy, line width=0.8pt] ([xshift=0.5cm, yshift=0.5cm]current page.south west) circle (2.5pt);
        \draw[bpNavy, line width=0.8pt] ([xshift=-0.5cm, yshift=0.5cm]current page.south east) circle (2.5pt);
        \draw[bpNavy, line width=0.8pt] ([xshift=0.5cm, yshift=-0.5cm]current page.north west) circle (2.5pt);
        \draw[bpNavy, line width=0.8pt] ([xshift=-0.5cm, yshift=-0.5cm]current page.north east) circle (2.5pt);

        % Banner superior
        \fill[bpNavy] ([xshift=0.5cm, yshift=-1.6cm]current page.north west) rectangle ([xshift=-0.5cm, yshift=-0.5cm]current page.north east);
        \node[anchor=center, text=white, font=\large\bfseries\sffamily] at ([yshift=-1.05cm]current page.north) {
            EL PLANO ARQUITECTÓNICO DE LA MATRIZ DE CONSISTENCIA
        };
        
        % Contenido Central
        \node[anchor=center, align=center] at ([yshift=0.2cm]current page.center) {
            {\color{bpNavy}\large\textbf{Arquitectura y Alineación Rigurosa de la Investigación}}\\[0.2cm]
            {\color{bpDark}\small Tesis: ``Diseño e implementación de una arquitectura SLM--RAG con QLoRA para la generación contextualizada en Quechua Collao''}\\[0.35cm]
            \begin{tcolorbox}[width=13.0cm, colback=bpLightGray, colframe=bpNavy!60!black, arc=1.5mm, center]
                \centering\small\sffamily
                \textbf{El ADN Inmutable:} $[\text{PROCESO: P}] + [\text{TECNOLOGÍA: T}] + [\text{CONTEXTO: C}]$\\
                \scriptsize Guía metodológica basada en las tres leyes de la investigación científica aplicada a la tesis
            \end{tcolorbox}
        };

        % Autores y cátedra
        \node[anchor=south, align=center, font=\tiny\sffamily, text=bpDark] at ([yshift=0.75cm]current page.south) {
            \textbf{Tesista:} Br. Edmil Jampier Saire Bustamante $\mid$ \textbf{Cátedra:} Msc. Raúl Huillca Huallparimachi (Seminario de Investigación II)\\
            \textbf{UNSAAC} -- Facultad de Ingeniería Eléctrica, Electrónica, Informática y Mecánica $\mid$ Cusco, 2026
        };
    \end{tikzpicture}
\end{frame}

% ==========================================================
% SLIDE 2: LAS 3 VARIABLES INMUTABLES (DICCIONARIO DE DATOS)
% ==========================================================
\begin{frame}{Las Tres Variables Inmutables de la Tesis}
    \vspace{-0.3cm}
    \begin{center}
    \scriptsize\textbf{Todo proyecto de investigación se reduce a tres variables inmutables:}
    \end{center}
    
    \vspace{0.1cm}
    \begin{tikzpicture}[node distance=0.25cm and 0.25cm, auto]
        % Variable P
        \node (varP) [draw=bpBlue, fill=bpBlue!8!white, rounded corners=1.5mm, line width=1.1pt, text width=3.8cm, minimum height=3.2cm, align=left] {
            \textbf{\color{bpBlue}\small [P] PROCESO}\\[0.05cm]
            \scriptsize\textbf{La acción metodológica}\\[0.15cm]
            \tiny\textbf{Definición en el Plan:}\\
            \color{bpDark}Diseño e implementación de una arquitectura de software \textbf{SLM--RAG offline} (sistema de procesamiento local autónomo sin conexión).
        };

        % Conector +
        \node (plus1) [right=0.15cm of varP, font=\Large\bfseries\color{bpNavy}] {+};

        % Variable T
        \node (varT) [right=0.15cm of plus1, draw=bpGold!90!black, fill=bpGold!10!white, rounded corners=1.5mm, line width=1.1pt, text width=3.8cm, minimum height=3.2cm, align=left] {
            \textbf{\color{bpGold!90!black}\small [T] TECNOLOGÍA}\\[0.05cm]
            \scriptsize\textbf{El núcleo técnico}\\[0.15cm]
            \tiny\textbf{Definición en el Plan:}\\
            \color{bpDark}Modelo compacto cuantizado a 4 bits con \textbf{adaptadores QLoRA} y recuperación local indexada.
        };

        % Conector +
        \node (plus2) [right=0.15cm of varT, font=\Large\bfseries\color{bpNavy}] {+};

        % Variable C
        \node (varC) [right=0.15cm of plus2, draw=bpRust, fill=bpRust!8!white, rounded corners=1.5mm, line width=1.1pt, text width=3.8cm, minimum height=3.2cm, align=left] {
            \textbf{\color{bpRust}\small [C] CONTEXTO}\\[0.05cm]
            \scriptsize\textbf{El ámbito / problema}\\[0.15cm]
            \tiny\textbf{Definición en el Plan:}\\
            \color{bpDark}Generación y anclaje factual de respuestas en \textbf{Quechua Collao} en entornos sin conectividad digital.
        };
    \end{tikzpicture}

    \vspace{0.2cm}
    \begin{tcolorbox}[enhanced, colback=bpLightGray, colframe=bpNavy, width=13.0cm, center, arc=1.2mm, top=1.2mm, bottom=1.2mm]
        \centering\tiny\sffamily
        \textbf{\color{bpNavy}Ley 1 (El Título es la Ley):} El título actúa como el genoma del proyecto. Los problemas, objetivos e hipótesis son solo \textbf{mutaciones gramaticales} de $[P]$, $[T]$ y $[C]$, sin agregar variables huérfanas.
    \end{tcolorbox}
\end{frame}

% ==========================================================
% SLIDE 3: DERIVACIÓN MACRO (NIVEL GENERAL)
% ==========================================================
\begin{frame}{Derivación Macro: Transformación Sintáctica del ADN}
    \vspace{-0.3cm}
    \begin{center}
    \tiny
    \begin{tabularx}{\textwidth}{l p{3.2cm} X}
        \toprule
        \textbf{\color{bpNavy}Elemento} & \textbf{\color{bpNavy}Estructura Gramatical} & \textbf{\color{bpNavy}Expresión en la Tesis} \\
        \midrule
        \textbf{Título Base} & 
        \textbf{\color{bpBlue}[P]} + \textbf{\color{bpGold!90!black}[T]} + \textbf{\color{bpRust}[C]} & 
        Diseño e implementación de una arquitectura SLM--RAG \textbf{\color{bpBlue}[P]} con QLoRA \textbf{\color{bpGold!90!black}[T]} para la generación contextualizada en Quechua Collao \textbf{\color{bpRust}[C]}. \\
        \addlinespace[0.15cm]
        \textbf{Problema General} & 
        ¿Cómo \textbf{\color{bpBlue}[P]} mediante \textbf{\color{bpGold!90!black}[T]} influye en \textbf{\color{bpRust}[C]}? & 
        ¿De qué manera el diseño e implementación de una arquitectura SLM--RAG \textbf{\color{bpBlue}[P]} con QLoRA \textbf{\color{bpGold!90!black}[T]} influye en la generación contextualizada en Quechua Collao \textbf{\color{bpRust}[C]}? \\
        \addlinespace[0.15cm]
        \textbf{Objetivo General} & 
        \textbf{\color{bpBlue}[Verbo P]} + \textbf{\color{bpGold!90!black}[T]} + \textbf{\color{bpRust}[C]} & 
        Diseñar e implementar una arquitectura SLM--RAG \textbf{\color{bpBlue}[P]} con QLoRA \textbf{\color{bpGold!90!black}[T]} para la generación contextualizada en Quechua Collao \textbf{\color{bpRust}[C]}. \\
        \addlinespace[0.15cm]
        \textbf{Hipótesis General} & 
        La aplicación de \textbf{\color{bpBlue}[P]} mediante \textbf{\color{bpGold!90!black}[T]} optimiza/garantiza \textbf{\color{bpRust}[C]} & 
        La implementación de la arquitectura SLM--RAG \textbf{\color{bpBlue}[P]} con QLoRA \textbf{\color{bpGold!90!black}[T]} optimiza la asimilación del contexto e incrementa la fidelidad factual en Quechua Collao \textbf{\color{bpRust}[C]}. \\
        \bottomrule
    \end{tabularx}
    \end{center}

    \vspace{0.15cm}
    \begin{tcolorbox}[colback=bpGreen!10!white, colframe=bpGreen!80!black, arc=1.2mm, left=2mm, right=2mm, top=1.2mm, bottom=1.2mm]
        \tiny\textbf{\color{bpGreen!80!black}Alineación Confirmada (100\% Consistente):}
        Sustantivo $\rightarrow$ Interrogación $\rightarrow$ Verbo Infinitivo $\rightarrow$ Afirmación Resolutiva. Ningún elemento conceptual nuevo ha sido inventado.
    \end{tcolorbox}
\end{frame}

% ==========================================================
% SLIDE 4: DESGLOSE CRONOLÓGICO (LOS 3 NODOS ESPECÍFICOS)
% ==========================================================
\begin{frame}{Desglose Específico: Secuencia Cronológica de Fases}
    \vspace{-0.3cm}
    \centering
    \scriptsize\textbf{Los Objetivos Específicos cronometran el proceso metodológico en 3 fases de ingeniería:}
    
    \vspace{0.25cm}
    \begin{tikzpicture}[node distance=0.25cm and 0.25cm, auto]
        % Nodo 1
        \node (node1) [draw=bpBlue, fill=bpBlue!8!white, rounded corners=1.5mm, line width=1.1pt, text width=4.0cm, minimum height=3.1cm, align=left] {
            \textbf{\color{bpBlue}\small FASE 1: Preparación}\\[0.05cm]
            \scriptsize\textbf{Node 1: Datos y Recuperación}\\[0.12cm]
            \tiny\textbf{Acción Metodológica:}\\
            \color{bpDark}Curar, estructurar e indexar semánticamente la base de conocimiento documental bilingüe para el motor de recuperación local en Quechua Collao.
        };

        % Flecha 1 -> 2
        \draw[-{Stealth[length=2.5mm]}, line width=1.2pt, color=bpNavy] ([xshift=0.08cm]node1.east) -- ++(0.35,0);

        % Nodo 2
        \node (node2) [right=0.45cm of node1, draw=bpGold!90!black, fill=bpGold!10!white, rounded corners=1.5mm, line width=1.1pt, text width=4.0cm, minimum height=3.1cm, align=left] {
            \textbf{\color{bpGold!90!black}\small FASE 2: Núcleo Técnico}\\[0.05cm]
            \scriptsize\textbf{Node 2: Adaptadores QLoRA}\\[0.12cm]
            \tiny\textbf{Acción Metodológica:}\\
            \color{bpDark}Implementar y acoplar adaptadores QLoRA en el SLM cuantizado a 4 bits para maximizar la asimilación del contexto RAG.
        };

        % Flecha 2 -> 3
        \draw[-{Stealth[length=2.5mm]}, line width=1.2pt, color=bpNavy] ([xshift=0.08cm]node2.east) -- ++(0.35,0);

        % Nodo 3
        \node (node3) [right=0.45cm of node2, draw=bpGreen!80!black, fill=bpGreen!8!white, rounded corners=1.5mm, line width=1.1pt, text width=4.0cm, minimum height=3.1cm, align=left] {
            \textbf{\color{bpGreen!80!black}\small FASE 3: Validación}\\[0.05cm]
            \scriptsize\textbf{Node 3: Benchmarking Computacional}\\[0.12cm]
            \tiny\textbf{Acción Metodológica:}\\
            \color{bpDark}Evaluar empíricamente el rendimiento (latencia/RAM), fidelidad léxica (chrF++, BLEU) y anclaje factual RAG (framework Ragas).
        };
    \end{tikzpicture}

    \vspace{0.2cm}
    \begin{tcolorbox}[colback=bpLightGray, colframe=bpBlue!70!black, width=13.5cm, center, arc=1.2mm, top=1.2mm, bottom=1.2mm]
        \centering\tiny\sffamily
        \textbf{\color{bpNavy}Ley 3 (Secuencia Lógica):} Los objetivos específicos constituyen los hitos técnicos cronológicos que resuelven secuencialmente la arquitectura propuesta de principio a fin.
    \end{tcolorbox}
\end{frame}

% ==========================================================
% SLIDE 5: ALINEACIÓN DETALLADA DE LOS 3 NODOS
% ==========================================================
\begin{frame}{Correspondencia Específica: Problemas vs. Objetivos}
    \vspace{-0.3cm}
    \begin{columns}[T]
        \begin{column}{0.32\textwidth}
            \begin{tcolorbox}[enhanced, colback=white, colframe=bpBlue, title=\textbf{\scriptsize Nodo 1: Datos y Recuperación}, arc=1.5mm, left=1.5mm, right=1.5mm, top=1.5mm, bottom=1.5mm]
                \tiny
                \textbf{\color{bpBlue}Problema Específico 1:}\\
                ¿Cómo estructurar e indexar semánticamente una base documental bilingüe para optimizar la recuperación local en Quechua Collao?\\[0.15cm]
                \textbf{\color{bpBlue}Objetivo Específico 1:}\\
                Curar, estructurar e indexar semánticamente una base documental bilingüe para optimizar la recuperación local en Quechua Collao.
            \end{tcolorbox}
        \end{column}
        
        \begin{column}{0.32\textwidth}
            \begin{tcolorbox}[enhanced, colback=white, colframe=bpGold!90!black, title=\textbf{\scriptsize Nodo 2: Núcleo Técnico}, arc=1.5mm, left=1.5mm, right=1.5mm, top=1.5mm, bottom=1.5mm]
                \tiny
                \textbf{\color{bpGold!90!black}Problema Específico 2:}\\
                ¿De qué manera los adaptadores QLoRA del SLM cuantizado optimizan la utilización de contexto bajo restricciones offline?\\[0.15cm]
                \textbf{\color{bpGold!90!black}Objetivo Específico 2:}\\
                Implementar y acoplar adaptadores QLoRA en el SLM cuantizado para maximizar la asimilación del contexto RAG en inferencia offline.
            \end{tcolorbox}
        \end{column}

        \begin{column}{0.32\textwidth}
            \begin{tcolorbox}[enhanced, colback=white, colframe=bpGreen!80!black, title=\textbf{\scriptsize Nodo 3: Validación Computacional}, arc=1.5mm, left=1.5mm, right=1.5mm, top=1.5mm, bottom=1.5mm]
                \tiny
                \textbf{\color{bpGreen!80!black}Problema Específico 3:}\\
                ¿Qué nivel de fidelidad léxica, consistencia factual y eficiencia alcanza la arquitectura frente al modelo base y RAG estándar?\\[0.15cm]
                \textbf{\color{bpGreen!80!black}Objetivo Específico 3:}\\
                Evaluar computacionalmente el rendimiento, fidelidad lingüística y anclaje factual frente al SLM base y RAG estándar.
            \end{tcolorbox}
        \end{column}
    \end{columns}

    \vspace{0.2cm}
    \begin{tcolorbox}[colback=bpLightGray, colframe=bpNavy, arc=1.2mm, left=2mm, right=2mm, top=1.2mm, bottom=1.2mm]
        \tiny\textbf{\color{bpNavy}Alineación 1 a 1:} Cada fase responde a una pregunta de ingeniería con su respectivo objetivo operativo verificable.
    \end{tcolorbox}
\end{frame}

% ==========================================================
% SLIDE 6: LA MATRIZ DE CONSISTENCIA COMPLETA
% ==========================================================
\begin{frame}{La Matriz de Consistencia Completa del Proyecto}
    \vspace{-0.35cm}
    \begin{center}
    \tiny
    \begin{tabularx}{\textwidth}{|>{\raggedright\arraybackslash}p{3.1cm}|>{\raggedright\arraybackslash}p{3.3cm}|>{\raggedright\arraybackslash}p{3.3cm}|>{\raggedright\arraybackslash}X|}
        \hline
        \rowcolor{bpNavy}
        \multicolumn{4}{|c|}{\textcolor{white}{\textbf{TÍTULO: Diseño e implementación de una arquitectura SLM--RAG con QLoRA para la generación contextualizada en Quechua Collao}}} \\
        \hline
        \rowcolor{bpLightGray}
        \textbf{\color{bpNavy}Problemas de Investigación} & \textbf{\color{bpNavy}Objetivos de la Tesis} & \textbf{\color{bpNavy}Hipótesis de Trabajo} & \textbf{\color{bpNavy}Variables e Instrumentos} \\
        \hline
        \multicolumn{4}{|c|}{\cellcolor{bpBlue!10}\textbf{\color{bpNavy}NIVEL GENERAL}} \\
        \hline
        ¿De qué manera el diseño e implementación de una arquitectura SLM--RAG con QLoRA influye en la generación contextualizada en Quechua Collao? & 
        Diseñar e implementar una arquitectura SLM--RAG con QLoRA para la generación contextualizada en Quechua Collao. & 
        La arquitectura SLM--RAG con QLoRA optimiza la asimilación del contexto e incrementa la fidelidad factual en Quechua Collao en inferencia offline. & 
        \textbf{VI:} Arquitectura SLM--RAG con QLoRA.\newline
        \textbf{VDs:} Rendimiento, calidad PLN y fidelidad RAG.\newline
        \textbf{Inst:} \texttt{llama.cpp}, Ragas, SacreBLEU. \\
        \hline
        \multicolumn{4}{|c|}{\cellcolor{bpBlue!10}\textbf{\color{bpNavy}NIVEL ESPECÍFICO (Fases Metodológicas de Ingeniería)}} \\
        \hline
        \textbf{PE1 (Datos):} ¿Cómo estructurar e indexar semánticamente una base documental bilingüe para optimizar la recuperación en Quechua Collao? & 
        \textbf{OE1:} Curar, estructurar e indexar semánticamente una base documental bilingüe para optimizar la recuperación local. & 
        \textbf{HE1:} La segmentación morfo-léxica e indexación semántica local maximizan la precisión de recuperación documental. & 
        \textbf{Var:} Precisión de indexación.\newline
        \textbf{Inst:} \texttt{sqlite-vec}, FAISS, \texttt{multilingual-e5-small}. \\
        \hline
        \textbf{PE2 (Núcleo):} ¿De qué manera los adaptadores QLoRA del SLM optimizan la utilización de contexto bajo restricciones offline? & 
        \textbf{OE2:} Implementar y acoplar adaptadores QLoRA en el SLM cuantizado para maximizar la asimilación del contexto RAG offline. & 
        \textbf{HE2:} La adaptación de bajo rango con QLoRA incrementa la fidelidad factual sin penalizar hardware modesto. & 
        \textbf{Var:} Latencia (ms), TPS, RAM (GB).\newline
        \textbf{Inst:} \texttt{llama.cpp}, \texttt{peft}, \texttt{psutil}. \\
        \hline
        \textbf{PE3 (Validación):} ¿Qué nivel de fidelidad léxica, consistencia factual y eficiencia alcanza la arquitectura frente al modelo base y RAG estándar? & 
        \textbf{OE3:} Evaluar computacionalmente el rendimiento, fidelidad lingüística y anclaje factual frente al SLM base y RAG estándar. & 
        \textbf{HE3:} La arquitectura con QLoRA supera significativamente al SLM base y RAG estándar en chrF++ y fidelidad Ragas. & 
        \textbf{Var:} chrF++, BLEU, Faithfulness.\newline
        \textbf{Inst:} Scripts Python, framework Ragas, Friedman/Wilcoxon. \\
        \hline
    \end{tabularx}
    \end{center}
\end{frame}

% ==========================================================
% SLIDE 7: LAS TRES LEYES Y CONTROL DE CALIDAD
% ==========================================================
\begin{frame}{Las Tres Leyes del Diseño y Panel de Control Final}
    \vspace{-0.3cm}
    \begin{columns}[T]
        \begin{column}{0.32\textwidth}
            \begin{tcolorbox}[enhanced, colback=bpLightGray, colframe=bpNavy, title=\textbf{\scriptsize 1. El Título es la Ley}, arc=1.5mm, left=1.5mm, right=1.5mm, top=1.5mm, bottom=1.5mm]
                \tiny
                \centering
                \textbf{\color{bpNavy}Consistencia de Variables}\\[0.1cm]
                \raggedright
                Cero variables huérfanas: problemas, objetivos e hipótesis derivan estrictamente de $[P]$, $[T]$ y $[C]$ sin variables pedagógicas extrañas.
            \end{tcolorbox}
        \end{column}
        
        \begin{column}{0.32\textwidth}
            \begin{tcolorbox}[enhanced, colback=bpLightGray, colframe=bpNavy, title=\textbf{\scriptsize 2. Transformación}, arc=1.5mm, left=1.5mm, right=1.5mm, top=1.5mm, bottom=1.5mm]
                \tiny
                \centering
                \textbf{\color{bpNavy}Mutación Gramatical}\\[0.1cm]
                \raggedright
                Los componentes generales son transformaciones sintácticas exactas (pregunta $\leftrightarrow$ verbo $\leftrightarrow$ afirmación), no ideas conceptuales distintas.
            \end{tcolorbox}
        \end{column}

        \begin{column}{0.32\textwidth}
            \begin{tcolorbox}[enhanced, colback=bpLightGray, colframe=bpNavy, title=\textbf{\scriptsize 3. Secuencia Lógica}, arc=1.5mm, left=1.5mm, right=1.5mm, top=1.5mm, bottom=1.5mm]
                \tiny
                \centering
                \textbf{\color{bpNavy}Cronología Ingenieril}\\[0.1cm]
                \raggedright
                Los objetivos específicos dividen el proceso en fases operativas de ingeniería DSR: \textbf{Datos} $\rightarrow$ \textbf{Adaptación con QLoRA} $\rightarrow$ \textbf{Evaluación}.
            \end{tcolorbox}
        \end{column}
    \end{columns}

    \vspace{0.15cm}
    \begin{tcolorbox}[enhanced, colback=white, colframe=bpGreen!80!black, title=\textbf{Panel de Control de Calidad y Auditoría Académica (100\% Aprobado)}, arc=1.5mm, left=2mm, right=2mm, top=1.5mm, bottom=1.5mm]
        \tiny
        \begin{tabular}{lll}
            \color{bpGreen}\bfseries [ $\checkmark$ ] & \textbf{Trazabilidad Total:} & Correspondencia matemática entre Título, PG, OG, HG y Fases Específicas. \\
            \color{bpGreen}\bfseries [ $\checkmark$ ] & \textbf{Cero Variables Huérfanas:} & Todas las variables específicas provienen del desglose de $[P]$ y $[T]$. \\
            \color{bpGreen}\bfseries [ $\checkmark$ ] & \textbf{Viabilidad Experimental:} & Cada objetivo específico cuenta con variables e instrumentos computacionales reproducibles. \\
            \color{bpGreen}\bfseries [ $\checkmark$ ] & \textbf{Rigor de Ingeniería:} & Enfoque DSR estructurado, autónomo y defendible ante el jurado de sistemas. \\
        \end{tabular}
    \end{tcolorbox}
\end{frame}

\end{document}
